\documentclass[11pt]{article}
\usepackage{mystyle}

\title{Rosen --- \S 8.5: Inclusion--Exclusion\\ \large Selected Exercises}
\author{Alston}
\date{Semester 2, 2026}

\begin{document}
\maketitle

% ======================================================================
\begin{exercise}{13}
    How many bit strings of length eight do not contain six
    consecutive 0s?
\end{exercise}

% ======================================================================
\begin{exercise}{14}
    How many permutations of the 26 letters of the English alphabet
    do not contain any of the strings \emph{fish}, \emph{rat}, or
    \emph{bird}?
\end{exercise}

% ======================================================================
\begin{exercise}{21}
    Write out the explicit formula given by the principle of
    inclusion--exclusion for the number of elements in the union of
    six sets when it is known that no three of these sets have a
    common intersection.
\end{exercise}

% ======================================================================
\begin{exercise}{22}
    Prove the principle of inclusion--exclusion using mathematical
    induction.
\end{exercise}

% ======================================================================
\begin{exercise}{29}
    Find a formula for the probability of the union of $n$ events in
    a sample space.
\end{exercise}

\end{document}